\documentclass[twocolumn]{iopjournal}

\usepackage{amsmath,amssymb,amsfonts}
\usepackage{amsthm}

\fancyhfoffset[L]{0mm}
\setlength{\marginparwidth}{0pt}
\setlength{\marginparsep}{0pt}

% ==========================================
% ZENODO PREPRINT HEADERS (ACTIVE)
% ==========================================
\renewcommand{\articletype}[1]{%
  {\vspace*{-8mm}\noindent \Large \sf Preprint}\par
  \vspace*{6mm}\noindent{\scriptsize \sf{\bfseries \MakeUppercase{#1}}}}

\fancyhead[L]{}
\fancyhead[R]{}

\hypersetup{pdftitle={Vortices of the Stellated Octahedron}, pdfauthor={Stephen Nelson}, pdfsubject={Gravity Field Dynamics (GFD)}, pdfkeywords={stellated octahedron, vortex, standing wave, tetrahedral symmetry, Field Origin, proof}}

\newtheorem{lemma}{Lemma}
\newtheorem{theorem}{Theorem}
\newtheorem{proposition}{Proposition}
\newtheorem{corollary}{Corollary}
\theoremstyle{definition}
\newtheorem{definition}{Definition}
\theoremstyle{remark}
\newtheorem*{remark}{Remark}

\begin{document}

\articletype{Proof}

\title{The Walsh Eigenmodes of the Cube, and What Field Closure Reads}

\author{Stephen Nelson}

\affil{Field Dynamics Research}

\ead{stephen@fielddynamics.org}

\begin{abstract}
The eight Walsh functions $1, x, y, z, yz, xz, xy, xyz$ on the corners of
the cube, the rows of the Hadamard matrix, are known to be its
eigenmodes. Under the tetrahedral symmetry $T_d$ of the stellated
octahedron they fall into two sectors: the uniform and parity modes $1,
xyz$ ($A_1$), and the six $T_2$ modes, three vector modes $x, y, z$ and
three quadrupole modes $yz, xz, xy$. We prove four things. Every Walsh
function is an eigenmode of a $T_d$-invariant coupling on the cube exactly
when TetA and TetB are alike; otherwise $1 \leftrightarrow xyz$ and
$x \leftrightarrow yz$, $y \leftrightarrow xz$, $z \leftrightarrow xy$ mix,
with angles in closed form, and in the full pixel the same pairing holds
mode by mode. The pixel's vortices are made of $T_2$ modes. Field closure,
with both roots of the tetrad active, is what reads the whole spectrum: the
linear readouts of one tetrahedron see only half of it, while those of both
see all eight modes, scalars reading $1$ and $xyz$ and the two frames the
six $T_2$ modes; frame closure, $\sum v_i = 0$, is forced by the symmetry
and is the same statement as the frame's blindness to the uniform mode. In
the archive's oscillator node, $\mu$ is the TetA share and depends on $1$
and $xyz$ alone, and the closure rule sees the rest only through its total
power. Finally, $\mu = x/(1+x)$ fixes the bonding $A_1$ mode: if the corner
values are field values it is proportional to $(g, a_0)$, and its chamber
parts are $g_N$ and $g - g_N$ of the field equation.
\end{abstract}

\keywords{Walsh functions, Hadamard, eigenmodes, standing wave,
tetrahedral symmetry, field closure, stellated octahedron, proof}

% ======================================================================

\section*{Setting}

\textbf{The cube.} The corners are $s \in \{\pm1\}^3$; TetA is the four
with $xyz = +1$ and TetB the four with $xyz = -1$~\cite{NelsonHadamard}.
The Walsh functions $w_S(s) = \prod_{d\in S}s_d$, $S \subseteq \{x, y,
z\}$, have \emph{grade} $\lvert S\rvert$ and are orthogonal, $w_S \cdot
w_T = 8\delta_{ST}$. On TetA, $yz = x$, $xz = y$, $xy = z$, and the Walsh
matrix is $H_2 \otimes \hat H$, with $H_2$ spatial parity and $\hat H$ the
bordered frame~\cite{NelsonHadamard}. As functions of direction the
grades have the angular characters s, p, d($t_{2g}$) and f($xyz$): grade 1
is a vector, odd under inversion, and grade 2 a quadratic form, even.

\textbf{Symmetries.} $F \cong \mathbb Z_2^3$ is the group of coordinate
sign flips (the cube's bit flips); $T_d$ (24 elements) the signed
permutations that keep TetA, and so TetB; $O_h = T_d \times \{\pm I\}$, the
inversion $-I$ swapping the chambers. A group acts on corner amplitudes by
$(g\psi)(gs) = \psi(s)$.

\textbf{The cube and the pixel.} An operator $H$ on the cube is a real
symmetric operator on the eight corner amplitudes: a Hamiltonian, or the
coupling matrix of eight oscillators, whose real eigenmodes are standing
waves. The \emph{pixel} is the cube with the octahedron and the Field
Origin, 15 sites, with the Hamiltonians of~\cite{NelsonVortex}.

\textbf{Two closures.} \emph{Frame closure} is Letter~1's
$v_1 + v_2 + v_3 + v_4 = 0$ for the four legs, the TetA
corners~\cite{NelsonL1}. \emph{Field closure} is Letter~3's: both roots of
the tetrad equation, TetA and TetB, active, nothing held at
zero~\cite{NelsonL3}. The oscillator node's update rule
(Section~5) is a third thing, which we call its \emph{closure rule}.

% ======================================================================

\section{The Walsh functions are the cube's eigenmodes}

\begin{theorem}[Known]
\label{thm:walsh}
$H$ commutes with every bit flip if and only if every Walsh function is
an eigenmode of $H$; such $H$ have 8 free parameters. In particular the
cube graph has $Aw_S = (3 - 2\lvert S\rvert)w_S$, and its Laplacian $L =
3I - A$ has $Lw_S = 2\lvert S\rvert\,w_S$~\cite{ODonnell2014}. Each $w_S$
is real, so as a mode of coupled oscillators it is a standing wave and
carries no current.
\end{theorem}

\begin{proof}
The $w_S$ are the eight distinct characters of $F$, so by Schur's
lemma~\cite{Serre1977} an operator commuting with $F$ is diagonal in that
basis, and conversely. A bit flip changes $w_S$ by $-1$ for each
coordinate in $S$, which gives $A$.
\end{proof}

% ======================================================================

\section{The pairing of the grades}

\begin{theorem}[The cube under $T_d$]
\label{thm:pairing}
\begin{enumerate}
\item[(i)] Under $T_d$ the corners carry $2A_1 \oplus 2T_2$: $1$ and
  $xyz$ are invariant and span $A_1$; the vector modes and the quadrupole
  modes are each a copy of $T_2$. Multiplying by $xyz$ commutes with $T_d$
  and maps grade 1 onto grade 2.
\item[(ii)] The most general $T_d$-invariant $H$ has six parameters:
  corner energies $\epsilon_A, \epsilon_B$; face-diagonal couplings
  $h_{AA}, h_{BB}$ inside each chamber; cube-edge and body-diagonal
  couplings $h_e, h_d$ between them. In the chamber basis ($u_{A,B}$ the
  normalised indicators, $f^d_{A,B}$ the coordinate $s_d$ restricted to a
  chamber and normalised),
  \begin{equation}
  H = \begin{pmatrix}\alpha_A & \beta\\ \beta & \alpha_B\end{pmatrix}
  \oplus
  \begin{pmatrix}\gamma_A & \delta\\ \delta & \gamma_B\end{pmatrix}^{\oplus 3},
  \label{eq:blocks}
  \end{equation}
  with $\alpha_{A,B} = \epsilon_{A,B} + 3h_{AA,BB}$, $\beta = 3h_e + h_d$,
  $\gamma_{A,B} = \epsilon_{A,B} - h_{AA,BB}$ and $\delta = h_e - h_d$;
  and $1 \propto u_A + u_B$, $xyz \propto u_A - u_B$, $x \propto f^x_A +
  f^x_B$, $yz \propto f^x_A - f^x_B$.
\item[(iii)] Every Walsh function is an eigenmode of $H$ if and only if
  $\epsilon_A = \epsilon_B$ and $h_{AA} = h_{BB}$, that is, if and only
  if $H$ is $O_h$-invariant (TetA and TetB alike). The eigenvalues are then
  $\alpha \pm \beta$ on $1, xyz$ and $\gamma \pm \delta$ on grades 1, 2.
\item[(iv)] Otherwise $H$ mixes exactly the pairs $1 \leftrightarrow xyz$
  and $x \leftrightarrow yz$, $y \leftrightarrow xz$, $z \leftrightarrow
  xy$. The eigenmodes are $(\cos\theta, \sin\theta)$ in the chamber basis
  of each block, with
  \begin{equation}
  \tan 2\theta_{A_1} = \frac{2\beta}{\alpha_A - \alpha_B},\quad
  \tan 2\theta_{T_2} = \frac{2\delta}{\gamma_A - \gamma_B},
  \label{eq:mixing}
  \end{equation}
  and weight $(\cos\theta + \sin\theta)^2/2$ on the lower grade of the pair.
\end{enumerate}
\end{theorem}

\begin{proof}
(i) By characters; $xyz$ is $T_d$-invariant, so multiplication by it is
equivariant. (ii) $T_d$ has two orbits of corners and four of corner pairs
(distances $2\sqrt2$ inside a chamber; $2$ and $2\sqrt3$ between). By
Schur, $H$ preserves the $A_1$ and $T_2$ parts and acts on $T_2$ as a
$2\times2$ matrix tensored with $I_3$; the entries follow from row sums and,
on $T_2$, from $\sum_{b \sim a}s_d(b) = s_d(a)$ along edges and $s_d(-a) =
-s_d(a)$. (iii) $(1, \pm1)$ are eigenvectors of a block iff its diagonal
entries agree, i.e.\ $\epsilon_A = \epsilon_B$ and $h_{AA} = h_{BB}$:
invariance under $-I$, which swaps the chambers and fixes the edge and
diagonal pairs. Under $O_h$ the grades are four inequivalent
representations. (iv) is a real symmetric $2\times2$ eigenproblem.
\end{proof}

This is the cube's form of a standard fact of tetrahedral chemistry: in
$T_d$, s and f$_{xyz}$ both carry $A_1$ and mix, and p and d($t_2$) both
carry $T_2$ and mix~\cite{Cotton1990}. The quadrupole modes behave as
directions only because $T_d$, having no inversion, cannot tell them from
vectors.

\begin{proposition}[The pixel]
\label{prop:pixel}
\begin{enumerate}
\item[(i)] If TetA and TetB are alike, the pixel has an eigenbasis whose
  corner parts are pure Walsh grades: grade 0 ($A_1$, even), 3 ($A_1$,
  odd), 1 ($T_2$, odd), 2 ($T_2$, even).
\item[(ii)] In general the corner part $c$ of an eigenmode at energy $E$
  satisfies $H_{\rm eff}(E)c = Ec$, with $H_{\rm eff}(E) = H_{cc} +
  H_{cr}(E - H_{rr})^{-1}H_{rc}$ ($r$: Field Origin and octahedron). This
  operator is $T_d$-invariant, so it has the form (\ref{eq:blocks}) with
  $E$-dependent entries: the pairing holds mode by mode, and each mode
  has its own angle.
\end{enumerate}
\end{proposition}

\begin{proof}
(i) The four projectors $P_{A_1}(1 \pm P_{-I})/2$ and $P_{T_2}(1 \pm
P_{-I})/2$ commute with every $O_h$-invariant $H$, and the corner part of
each lies in one grade (checked exactly). (ii) Eliminate the 7 non-corner
sites; the blocks are $T_d$-equivariant. At one parameter point the three
$T_2$ levels have grade-1 shares 0.83, 0.007 and 0.98 (numerical).
\end{proof}

% ======================================================================

\section{The vortices are made of $T_2$ modes}

\begin{corollary}
The pixel carries $4A_1 \oplus E \oplus 3T_2$ (two $T_2$ copies on the
corners, one on the octahedron). The rotoreflection $S$ has eigenvalue
$+i$ only on $T_2$, so the winding space $V_{+i}$~\cite{NelsonVortex}
takes one vector from each $T_2$ copy. Every vortex's corner part is a
combination of the six $T_2$ modes, and its octahedron part lies in the
octahedron's $T_2$ copy; the uniform and parity modes never wind.
\end{corollary}

\begin{proof}
$S$ acts as $1$ on $A_1$, with eigenvalues $\pm1$ on $E$, and with $i, -i,
-1$ on each $T_2$.
\end{proof}

\begin{proposition}[Two standing waves make a vortex]
\label{prop:cube}
On the cube graph, $H = tA$ with $t > 0$, $\psi = x + iy$ is an eigenmode at
the triply degenerate grade-1 energy $t$ (spanned by $x, y, z$). It winds
under the quarter turn about $\hat z$, $C_4\psi = -i\psi$, and carries the
current $-8t$ ($-t/2$ per unit norm) through every half-plane bounded by
the $z$ axis, counted as the physical flow $J_{i\to j} =
2\,\mathrm{Im}(\bar\psi_jH_{ji}\psi_i)$ in the right-handed
sense~\cite{NelsonVortex}. The standing waves $x$ and $y$ alone carry none.
\end{proposition}

\begin{proof}
$Ax = x$, $Ay = y$; $C_4$ sends $x \mapsto y$, $y \mapsto -x$, so $C_4\psi =
-i\psi$; the current is computed exactly through three half-planes.
\end{proof}

% ======================================================================

\section{What frame closure and field closure read}

Write $S_{A,B}$ for the sum of the amplitudes over a chamber and $R_{A,B}$
for its frame readout, $\psi \mapsto \sum_{s\in{\rm chamber}}\psi(s)\,s$.

\begin{theorem}[Readouts]
\label{thm:closure}
\begin{enumerate}
\item[(i)] Every $T_d$-equivariant linear map from TetA's four amplitudes
  to $\mathbb R^3$ is a multiple of the frame readout $R_A$.
\item[(ii)] Frame closure is forced by the symmetry: the legs are one
  $T_d$-orbit and $T_d$ fixes no nonzero vector, so $\sum v_i = 0$. It is
  the same equation as $R_A(1) = 0$, the frame's blindness to the uniform
  mode; $R_A$ is injective on the sum-zero modes.
\item[(iii)] One root: TetA's readouts $S_A, R_A$ have rank 4 on the eight
  corners and cannot tell $w_S$ from $w_{\bar S}$ ($\bar S$ the
  complement), the collapse on one tetrahedron~\cite{NelsonHadamard}: they
  read half of the spectrum.
\item[(iv)] Field closure: $(S_A, S_B, R_A, R_B)$ is an invertible readout
  of all eight corners. $S_A, S_B$ read exactly $1$ and $xyz$; $R_A, R_B$
  read exactly the six $T_2$ modes. The $T_d$-equivariant linear readouts
  are spanned by these, two scalars and two frames, so every one of them
  factors through the field-closed readout.
\end{enumerate}
\end{theorem}

\begin{proof}
TetA's corners carry $A_1 \oplus T_2$ and $\mathbb R^3$ carries $T_2$, so
$\mathrm{Hom}_{T_d}$ is one-dimensional (Schur). The $T_d$-orbit of $v_1$ is
the four legs; the stacked $g - I$ have rank 3. On eight corners
$\mathrm{Hom}_{T_d}(2A_1 \oplus 2T_2, T_2)$ and $\mathrm{Hom}_{T_d}(2A_1
\oplus 2T_2, A_1)$ are two-dimensional, and $R_{A,B}$, $S_{A,B}$ are
equivariant and independent. All four statements are also solved directly
from the linear equations.
\end{proof}

So a readout that uses one root reads half of the Walsh spectrum, and
field closure is what makes the reading complete. $\mu = S_A/(S_A + S_B)$
cannot be formed from TetA alone: an amplitude supported on TetA and the
uniform amplitude have the same TetA readouts and $\mu = 1$, $\frac12$.
The division into two sectors is a property of \emph{linear} readouts: a
quadratic invariant, such as $\lvert R_A\psi\rvert^2$, is a scalar that
reads the $T_2$ modes.

% ======================================================================

\section{What the oscillator node's closure rule reads}

The archive's octahedral resonance node~\cite{NelsonPixelated} drives eight
corner oscillators, mixes them, and forms a vertex state $v_k$ with Walsh
coefficients $c_S = \frac18\sum_kw_S(s_k)v_k$. It computes $\mu = (1 +
c_{xyz}/c_1)/2$, and its closure rule makes three updates at rate $\kappa =
0.02$: amplitudes toward $a^* = r/(1+r)$ (bounded rule) or $\min(r, 5)$
(original), with $r^2$ the mean energy of the two mixes; outward couplings
toward 1 or $\min(r, 5)$; inward couplings by $c_k \leftarrow c_k +
\kappa\,\mathrm{sign}(f_kb_k)(1 - \lvert c_k\rvert)$.

\begin{theorem}[The node]
\label{thm:node}
\begin{enumerate}
\item[(i)] $\mu = S_A/(S_A + S_B)$, the TetA share of the vertex state; it
  depends only on $1$ and $xyz$, and adding any $T_2$ modes leaves it
  unchanged.
\item[(ii)] The energy that drives the rule is $\sum_kv_k^2 = 8\sum_Sc_S^2$
  (Parseval), the total Walsh power. Amplitudes and outward couplings each
  have one target for all corners, so while no clamp acts the spread
  between corners contracts by exactly $1 - \kappa$ per step, and at equal
  amplitudes a step rescales the carrier's whole spectrum by one factor.
\item[(iii)] The inward-coupling update is corner-local and depends on the
  signal; it is inert in the archive only because every $c_k$ starts at its
  fixed point 1.
\end{enumerate}
\end{theorem}

\begin{proof}
(i) $c_1 = (S_A + S_B)/8$, $c_{xyz} = (S_A - S_B)/8$. (ii) Parseval; $a_k
- a_j \leftarrow (1 - \kappa)(a_k - a_j)$ under a common target; the
carrier is linear in a common amplitude. (iii) At $c_k = 1$ the increment
vanishes; started at 0.5 the couplings spread across corners while
relaxing to 1. All three are checked on the node's own code, 800 steps in
both modes.
\end{proof}

% ======================================================================

\section{$\mu$ and the mixing angle}

\begin{proposition}
\label{prop:mu}
Let $(\cos\theta, \sin\theta)$, $0 < \theta < \pi/2$, be the bonding mode
of the $A_1$ block (\ref{eq:blocks}). Its TetA share of amplitude is $1/(1
+ \tan\theta)$; the antibonding mode's, $1/(1 - \cot\theta)$, lies outside
$[0, 1]$. With $x = g/a_0$:
\begin{enumerate}
\item[(i)] \emph{Linear reading} (corner values are field values): $\mu =
  x/(1+x)$ if and only if $\tan\theta = a_0/g$, the mode proportional to
  $(g, a_0)$; equivalently $\beta/(\alpha_A - \alpha_B) = x/(x^2-1)$.
\item[(ii)] In that reading the mode's chamber parts, scaled to total $g$,
  are $\mu g = g^2/(g + a_0)$ and $(1-\mu)g = ga_0/(g + a_0)$: the TetA part
  is $g_N$ of the field equation $g_N = g^2/(g + a_0)$.
\item[(iii)] \emph{Quadratic reading} (the field is the intensity): $\mu =
  x/(1+x)$ if and only if $\tan\theta = \sqrt{a_0/g}$, the mode
  proportional to $(\sqrt g, \sqrt{a_0})$, with $\beta/(\alpha_A -
  \alpha_B) = \sqrt x/(x-1)$: the bond-orbital dictionary $V_3/V_2 = (x -
  1)/(2\sqrt x)$ of the diamond lattice.
\item[(iv)] The readings agree only at $x = 1$ ($\theta = \pi/4$, $\mu =
  \frac12$), where TetA and TetB are alike.
\end{enumerate}
\end{proposition}

\begin{proof}
The mode has chamber sums proportional to $\cos\theta$ and $\sin\theta$,
and chamber weights $\cos^2\theta$ and $\sin^2\theta$; the rest is algebra
with $\tan 2\theta = 2\tan\theta/(1 - \tan^2\theta)$ and (\ref{eq:mixing}).
\end{proof}

\begin{remark}[Which reading]
For the cube the linear reading is the consistent one. Every readout in
Sections~4 and~5 is linear in the corner values, and Letter~3 divides the
field itself, TetA carrying $\mu$ of $g$ and TetB $1 - \mu$, with general
relativity seeing $\mu g = g_N$~\cite{NelsonL3}; Proposition~\ref{prop:mu}(ii)
reproduces that division and the field equation exactly. The quadratic
reading is the right one when the field is the intensity of a lattice
wave, as in the diamond lattice's bond orbital. Which applies is a
statement about what the mode's amplitude represents, not about $\mu$.
\end{remark}

% ======================================================================

\section{Remarks}

\begin{remark}[What has been proved]
The cube holds the Walsh standing waves exactly when TetA and TetB are
alike, and in rotated pairs when they differ; in the pixel the pairing
holds mode by mode. The pixel's vortices are made of $T_2$ modes. Among
linear symmetric readouts, one root reads half of the spectrum, and field
closure reads all of it, in two sectors: the scalars read $1$ and $xyz$
and give $\mu$; the frames read the six $T_2$ modes. The node's closure
rule sees the rest only through the total power.
\end{remark}

\begin{remark}[Not claimed]
That the physical field is this node or reads linearly, or that any model's
bonding $A_1$ mode satisfies $\tan\theta = a_0/g$. Proposition~\ref{prop:mu}
makes that the test: compute the bonding mode of a lattice model and
compare its angle with $a_0/g$.
\end{remark}

% ======================================================================

\section*{Verification}

\texttt{verify\_walsh\_eigenmodes.py} runs 45 checks in six parts: the
cube's eigenmodes; the pairing on the cube, in all parameters, and in the
pixel (exactly for alike chambers, numerically mode by mode otherwise); the
vortices and the cube vortex; the readouts, solved from their equivariance
equations; the node, symbolically and on its own code
(\texttt{octahedral\_node.py} in~\cite{NelsonPixelated}); and the two
readings of $\mu$.

\begin{thebibliography}{9}

\bibitem{NelsonHadamard}
S.~Nelson,
\textit{The Stellated Octahedron, the Cube and the Hadamard Matrix},
Field Dynamics Research (2026),
\url{proofs/hadamard-cube-frame}.

\bibitem{NelsonVortex}
S.~Nelson,
\textit{Vortices of the Stellated Octahedron},
Field Dynamics Research (2026),
\url{proofs/vortices-of-the-stellated-octahedron}.

\bibitem{NelsonL1}
S.~Nelson,
\textit{The Coordinate-Space Geometry of the Tetrad Frame} (Letter~1),
Field Dynamics Research (2026).

\bibitem{NelsonL3}
S.~Nelson,
\textit{Gravity Field Dynamics: Field Closure of the Tetrad Frame}
(Letter~3), Field Dynamics Research (2026).

\bibitem{NelsonPixelated}
S.~Nelson,
\textit{The Universe Is Pixelated}, research note with the octahedral
resonance node, Field Dynamics Research (2026),
\url{research/universe-is-pixelated}.

\bibitem{ODonnell2014}
R.~O'Donnell,
\textit{Analysis of Boolean Functions}
(Cambridge University Press, Cambridge, 2014).

\bibitem{Serre1977}
J.-P.~Serre,
\textit{Linear Representations of Finite Groups}
(Springer, New York, 1977).

\bibitem{Cotton1990}
F.~A.~Cotton,
\textit{Chemical Applications of Group Theory}, 3rd ed.
(Wiley, New York, 1990).

\end{thebibliography}

\end{document}
